\makeatletter
\def\input@path{{./}{../}{../../}{../../../}{../../../../}{preamble/}{../preamble/}{../../preamble/}{../../../preamble/}}
\makeatother
\input{preamble_exams.tex}

\title{Grade 11 -- Term 1 \\ C9 Practice Activity with Solutions}
\author{}
\date{}

\begin{document}
\maketitle
\setcounter{tocdepth}{1}
\setcounter{secnumdepth}{2}
\phantomsection
\label{toc}
\tableofcontents

\section*{Evaluated criterion}
\begin{itemize}
\item[\textbf{C9.}] Calculates and compares critical values and probability areas from the standard normal and Student's $t$-distributions, selecting and justifying the appropriate distribution according to whether the population standard deviation is known or estimated from the sample.
\end{itemize}
\vspace{-4mm}

\section{Introduction}
\label{sec:c9-introduction}
\SectionProblemsTableOfContents{
\SectionProblemLink{sec:c9-choice}{Choosing the distribution} \\
\SectionProblemLink{sec:c9-areas}{Critical values and probability areas} \\
\SectionProblemLink{sec:c9-tables}{Tables for this activity}
}
\vspace{5mm}

\subsection{Choosing the distribution}
\label{sec:c9-choice}
The standard normal distribution, $Z\sim N(0,1)$, and Student's distribution, $T\sim t_\nu$, are continuous, symmetric about zero, and bell-shaped. Each has total area $1$. The standard normal has mean $0$ and standard deviation $1$. Student's distribution is a family indexed by its degrees of freedom $\nu$; its standard deviation is not generally $1$.

In this activity, all samples consist of independent random observations from a normal population. These assumptions make the following sample-mean models exact:
\[
\begin{aligned}
\text{Known population standard deviation }\sigma:
&\qquad Z=\frac{\bar X-\mu}{\sigma/\sqrt n}\sim N(0,1),\\
\text{Unknown }\sigma\text{, estimated using }S:
&\qquad T=\frac{\bar X-\mu}{S/\sqrt n}\sim t_{n-1}.
\end{aligned}
\]
Here $\mu$ is the population mean, $\bar X$ the random sample mean, and $S$ the random sample standard deviation; an observed sample standard deviation is written $s$. No calculation from raw data is needed. The cutoffs in this activity are dimensionless standardized values.

Estimating the mean leaves $n-1$ independent deviations, giving $\nu=n-1$. Replacing $\sigma$ by $s$ adds uncertainty, represented by heavier $t$-tails. For the same small tail probability, the positive $t$-cutoff is larger than the $z$-cutoff. As $\nu$ increases, $t_\nu$ approaches $N(0,1)$. A large sample does not make an unknown $\sigma$ known: use the appropriate $t$ row whenever $s$ estimates $\sigma$ here.

\begin{center}
\begin{tabularx}{0.96\textwidth}{@{}lXX@{}}
\toprule
Feature & Standard normal & Student's $t$\\
\midrule
Variability information & Population $\sigma$ known & Population $\sigma$ unknown; use sample $s$\\
Sample-mean denominator & $\sigma/\sqrt n$ & $s/\sqrt n$\\
Degrees of freedom & No row depending on $n$ & $\nu=n-1$\\
Shape and tails & One fixed curve & Heavier tails; approach normal as $\nu$ increases\\
\bottomrule
\end{tabularx}
\end{center}
\secInicio[sec:c9-choice][sec:c9-introduction]

\subsection{Critical values and probability areas}
\label{sec:c9-areas}
Write $W$ for whichever standardized variable is appropriate. A critical value is a cutoff leaving a specified probability area. For $c>0$, write $q=P(W>c)$. Then
\[
\begin{aligned}
P(W<-c)&=q,& P(W<c)&=1-q,\\
P(|W|>c)&=2q,& P(-c<W<c)&=1-2q.
\end{aligned}
\]
Define positive critical values by $P(Z>z_q)=q$ and $P(T>t_{\nu,q})=q$. A left-tail area $q<0.5$ requires the negative cutoff $-z_q$ or $-t_{\nu,q}$. For combined two-tail area $\alpha$, use $q=\alpha/2$; for central area $p$, use $q=(1-p)/2$. For an asymmetric interval, subtract both excluded tails from $1$. Endpoints do not affect probabilities for continuous distributions.
\secInicio[sec:c9-areas][sec:c9-introduction]

\subsection{Tables for this activity}
\label{sec:c9-tables}
Use only the tables below. Critical values are rounded to three decimal places and cumulative normal areas to four. Treat probabilities at printed critical values as table-based approximations. For cutoffs strictly between critical entries, report bounds without interpolation. Every required degree-of-freedom row and normal lookup is included.

\subsubsection*{Standard normal cumulative areas (selected entries)}
The entries are $\Phi(z)=P(Z<z)$ for positive $z$. Use
$\Phi(-z)=1-\Phi(z)$ and $P(Z>z)=1-\Phi(z)$.
\begin{center}
\begin{tabular}{r|r@{\qquad\qquad}r|r}
\toprule
$z$ & $\Phi(z)$ & $z$ & $\Phi(z)$\\
\midrule
0.000 & 0.5000 & 1.900 & 0.9713\\
1.000 & 0.8413 & 1.960 & 0.9750\\
1.282 & 0.9001 & 2.000 & 0.9772\\
1.500 & 0.9332 & 2.300 & 0.9893\\
1.645 & 0.9500 & 2.326 & 0.9900\\
1.833 & 0.9666 & 2.576 & 0.9950\\
\bottomrule
\end{tabular}
\end{center}

\subsubsection*{Positive upper-tail critical values}
In both tables below, $q$ denotes \textbf{one upper tail}. The second and third header rows convert it to a combined two-tail area or central area.
\begin{center}
\begingroup
\setlength{\tabcolsep}{8pt}
\begin{tabular}{c|rrrrr}
\toprule
Upper tail $q$ & $0.10$ & $0.05$ & $0.025$ & $0.01$ & $0.005$\\
Two tails $2q$ & $0.20$ & $0.10$ & $0.05$ & $0.02$ & $0.01$\\
Central $1-2q$ & $0.80$ & $0.90$ & $0.95$ & $0.98$ & $0.99$\\
\midrule
Normal $z_q$ & 1.282 & 1.645 & 1.960 & 2.326 & 2.576\\
\bottomrule
\end{tabular}

\medskip
\begin{tabular}{c|rrrrr}
\toprule
Upper tail $q$ & $0.10$ & $0.05$ & $0.025$ & $0.01$ & $0.005$\\
Two tails $2q$ & $0.20$ & $0.10$ & $0.05$ & $0.02$ & $0.01$\\
Central $1-2q$ & $0.80$ & $0.90$ & $0.95$ & $0.98$ & $0.99$\\
\midrule
Degrees of freedom $\nu$ & \multicolumn{5}{c}{Student's positive critical value $t_{\nu,q}$}\\
\midrule
5 & 1.476 & 2.015 & 2.571 & 3.365 & 4.032\\
9 & 1.383 & 1.833 & 2.262 & 2.821 & 3.250\\
11 & 1.363 & 1.796 & 2.201 & 2.718 & 3.106\\
14 & 1.345 & 1.761 & 2.145 & 2.624 & 2.977\\
19 & 1.328 & 1.729 & 2.093 & 2.539 & 2.861\\
24 & 1.318 & 1.711 & 2.064 & 2.492 & 2.797\\
29 & 1.311 & 1.699 & 2.045 & 2.462 & 2.756\\
59 & 1.296 & 1.671 & 2.001 & 2.391 & 2.662\\
119 & 1.289 & 1.658 & 1.980 & 2.358 & 2.618\\
\bottomrule
\end{tabular}
\endgroup
\end{center}
Larger positive cutoffs leave smaller upper-tail areas. For example, in row $9$, $1.833<2.00<2.262$ implies $0.025<P(T>2.00)<0.05$. The displayed precision is sufficient for all bounds requested here. The cumulative and critical normal tables may differ slightly because of rounding; use the critical table for requested cutoffs and the cumulative table for areas at given normal values.
\secInicio[sec:c9-tables][sec:c9-introduction]

\section{Comparing \texorpdfstring{$Z$ and $t$}{Z and t}}
\label{sec:c9-comparing}
\SectionProblemsTableOfContents{
\SectionProblemLink{sec:c9-critical}{One tail, two tails, and central cutoffs} \\
\SectionProblemLink{sec:c9-fixed}{Areas at the same cutoff} \\
\SectionProblemLink{sec:c9-convergence}{Increasing degrees of freedom}
}
\vspace{5mm}
Attempt each task before reading its worked solution.

\subsection{One tail, two tails, and central cutoffs}
\label{sec:c9-critical}
Compare $Z\sim N(0,1)$ and $T\sim t_9$.
\QuestionsTasks[sec:c9-comparing]{{9}/{sol:c9-critical}/{Find (a) a cutoff with upper-tail area $0.05$; (b) both cutoffs with combined two-tail area $0.05$; and (c) both cutoffs enclosing the middle $80\%$. Compare the two models in each case.}}
\subsubsection*{C9 Solution}
\phantomsection\label{sol:c9-critical}
\begin{center}
\begin{tabular}{lccc}
\toprule
Request & Upper tail $q$ & Normal cutoff(s) & $t_9$ cutoff(s)\\
\midrule
(a) One upper tail & 0.05 & 1.645 & 1.833\\
(b) Two tails & 0.025 & $\pm1.960$ & $\pm2.262$\\
(c) Central area & 0.10 & $\pm1.282$ & $\pm1.383$\\
\bottomrule
\end{tabular}
\end{center}
The positive $t$-cutoff exceeds the corresponding $z$-cutoff by $0.188$, $0.302$, and $0.101$, respectively. Heavier $t$-tails require wider symmetric ranges to enclose the same central probability.
\secInicio[sec:c9-critical][sec:c9-comparing]

\subsection{Areas at the same cutoff}
\label{sec:c9-fixed}
Compare the models at the fixed cutoff $1.833$.
\QuestionsTasks[sec:c9-comparing]{{9}/{sol:c9-fixed}/{Calculate the upper-tail, combined two-tail, and central probabilities for $Z$ and $T\sim t_9$ at symmetric cutoffs $\pm1.833$. Then bound the upper-tail and central probabilities for $T\sim t_9$ at $\pm2.00$, and compare with the normal areas.}}
\subsubsection*{C9 Solution}
For $Z$, $\Phi(1.833)=0.9666$; for $t_9$, $1.833$ is in column $q=0.05$.
\phantomsection\label{sol:c9-fixed}
\[
\begin{array}{c|ccc}
\text{Model} & P(W>1.833) & P(|W|>1.833) & P(-1.833<W<1.833)\\ \hline
Z & 0.0334 & 0.0668 & 0.9332\\
t_9 & 0.0500 & 0.1000 & 0.9000
\end{array}
\]
At $2.00$, $1.833<2.00<2.262$ gives
\[
0.025<P(T>2.00)<0.05,\qquad
0.90<P(-2.00<T<2.00)<0.95.
\]
For $Z$, the upper tail is $1-0.9772=0.0228$ and the central area is $1-2(0.0228)=0.9544$. Thus the $t_9$ upper tail is larger and its central area is smaller even using only these bounds.
\secInicio[sec:c9-fixed][sec:c9-comparing]

\subsection{Increasing degrees of freedom}
\label{sec:c9-convergence}
\QuestionsTasks[sec:c9-comparing]{{9}/{sol:c9-convergence}/{For $\nu=5,19,59,119$, find the positive cutoff enclosing central area $0.95$ together with its negative counterpart. Calculate each positive cutoff's difference from $z_{0.025}$. Explain the trend.}}
\subsubsection*{C9 Solution}
\phantomsection\label{sol:c9-convergence}
Central area $0.95$ gives $q=0.025$ and $z_{0.025}\approx1.960$.
\begin{center}
\begin{tabular}{c|rr}
\toprule
$\nu$ & $t_{\nu,0.025}$ & $t_{\nu,0.025}-1.960$\\
\midrule
5 & 2.571 & 0.611\\
19 & 2.093 & 0.133\\
59 & 2.001 & 0.041\\
119 & 1.980 & 0.020\\
\bottomrule
\end{tabular}
\end{center}
The symmetric cutoffs are the positive and negative values in each row. The decreasing difference shows convergence toward the normal cutoffs, while the finite-$\nu$ values remain larger.
\secInicio[sec:c9-convergence][sec:c9-comparing]

\section{Contextual Practice}
\label{sec:c9-contextual}
For every problem, select the distribution and justify it using the status of $\sigma$. Calculate $\nu$ when needed, classify the area, and show the table row and column or cumulative entry used. All sample-mean models satisfy the assumptions in the introduction.
\SectionProblemsTableOfContents{
\SectionProblemLink{sec:c9-households}{Household expenditure} \\
\SectionProblemLink{sec:c9-investment}{Investment returns} \\
\SectionProblemLink{sec:c9-deliveries}{Delivery times} \\
\SectionProblemLink{sec:c9-production}{Production costs} \\
\SectionProblemLink{sec:c9-revenue}{Business revenues} \\
\SectionProblemLink{sec:c9-wages}{Wage studies}
}
\vspace{5mm}

\subsection{Household expenditure}
\label{sec:c9-households}
A survey uses $n=25$ monthly household expenditures. The population standard deviation is known to be $\sigma=120$ monetary units. Let $W$ be the standardized sample mean.
\QuestionsTasks[sec:c9-contextual]{{9}/{sol:c9-households}/{Find the cutoff $c$ with $P(W<c)=0.01$ and the symmetric cutoffs enclosing the middle $95\%$.}}
\subsubsection*{C9 Solution}
\phantomsection\label{sol:c9-households}
Use $Z$ because population $\sigma$ is known; no degrees-of-freedom row is needed. For the left tail, column $q=0.01$ gives $c\approx-2.326$. For the central area, $q=0.025$ gives cutoffs $\pm1.960$.
\secInicio[sec:c9-households][sec:c9-contextual]

\subsection{Investment returns}
\label{sec:c9-investment}
An analyst samples $n=10$ investment returns. Population $\sigma$ is unknown and is estimated by $s=1.8$ percentage points. Let $W$ be the standardized sample mean.
\QuestionsTasks[sec:c9-contextual]{{9}/{sol:c9-investment}/{Find a positive cutoff exceeded with probability $0.05$. Calculate $P(W<-1.833)$ and $P(-1.833<W<2.262)$.}}
\subsubsection*{C9 Solution}
\phantomsection\label{sol:c9-investment}
Use $t$ because $s$ estimates unknown $\sigma$; $\nu=10-1=9$. Row $9$, column $0.05$, gives $c\approx1.833$. Symmetry gives $P(W<-1.833)\approx0.05$. The asymmetric interval excludes tails $0.05$ and $0.025$, so its area is approximately $1-0.05-0.025=0.925$.
\secInicio[sec:c9-investment][sec:c9-contextual]

\subsection{Delivery times}
\label{sec:c9-deliveries}
A logistics company samples $n=15$ deliveries. The population standard deviation is unknown; the observed sample has $s=6$ minutes. Let $W$ be the standardized sample mean.
\QuestionsTasks[sec:c9-contextual]{{9}/{sol:c9-deliveries}/{Find symmetric cutoffs with a total outside area of $0.10$. Then bound $P(W<-2.30)$ and $P(W>-2.30)$.}}
\subsubsection*{C9 Solution}
\phantomsection\label{sol:c9-deliveries}
Use $t_{14}$ because $\sigma$ is estimated using $s$ and $\nu=15-1=14$. The two-tail request gives $q=0.05$, hence cutoffs $\pm1.761$. Since $2.145<2.30<2.624$, the upper-tail area lies between $0.01$ and $0.025$. By symmetry and complementation,
\[
0.01<P(W<-2.30)<0.025,\qquad
0.975<P(W>-2.30)<0.99.
\]
\secInicio[sec:c9-deliveries][sec:c9-contextual]

\subsection{Production costs}
\label{sec:c9-production}
A manufacturer samples $n=20$ batch costs from a population with known $\sigma=40$ monetary units. Let $W$ be the standardized sample mean.
\QuestionsTasks[sec:c9-contextual]{{9}/{sol:c9-production}/{Calculate $P(W>1.50)$, $P(|W|>1.50)$, and $P(-1.50<W<1.50)$. Find the symmetric cutoffs enclosing the middle $98\%$.}}
\subsubsection*{C9 Solution}
\phantomsection\label{sol:c9-production}
Use $Z$ because $\sigma$ is known. From $\Phi(1.50)=0.9332$, the areas are $0.0668$, $0.1336$, and $0.8664$, respectively. Central area $0.98$ gives $q=0.01$, hence cutoffs $\pm2.326$.
\secInicio[sec:c9-production][sec:c9-contextual]

\subsection{Business revenues}
\label{sec:c9-revenue}
A study samples $n=30$ businesses and estimates unknown population $\sigma$ using $s=8$ thousand monetary units. Let $W$ be the standardized sample mean.
\QuestionsTasks[sec:c9-contextual]{{9}/{sol:c9-revenue}/{Bound the upper-tail, combined two-tail, and central areas at symmetric cutoffs $\pm1.90$. Compare these areas with a second study whose population $\sigma$ is known, using the same standardized cutoffs.}}
\subsubsection*{C9 Solution}
\phantomsection\label{sol:c9-revenue}
The first study uses $t_{29}$: $\sigma$ is unknown and $\nu=30-1=29$. Since $1.699<1.90<2.045$,
\[
0.025<P(W>1.90)<0.05,\quad
0.05<P(|W|>1.90)<0.10,\quad
0.90<P(-1.90<W<1.90)<0.95.
\]
The second study uses $Z$. Its corresponding areas are $0.0287$, $0.0574$, and $0.9426$, from $\Phi(1.90)=0.9713$. These normal areas lie within the coarse $t$ bounds, so the bounds alone do not quantify their differences. The heavier-tail property tells us the $t$ tail areas are larger and its central area smaller.
\secInicio[sec:c9-revenue][sec:c9-contextual]

\subsection{Wage studies}
\label{sec:c9-wages}
Study A samples $n=6$ wages with known population $\sigma=200$ monetary units. Study B samples $n=120$ wages with unknown $\sigma$, estimated by $s=200$ monetary units.
\QuestionsTasks[sec:c9-contextual]{{9}/{sol:c9-wages}/{For each study, find the negative cutoff leaving $0.025$ below it and the symmetric cutoffs enclosing the middle $90\%$. Explain why sample size and the equal numerical standard deviations do not determine the distribution choice.}}
\subsubsection*{C9 Solution}
\phantomsection\label{sol:c9-wages}
A uses $Z$ because $\sigma$ is known, even with $n=6$: the left cutoff is $-1.960$ and the central cutoffs are $\pm1.645$. B uses $t_{119}$ because its standard deviation is an estimate: the left cutoff is $-1.980$ and the central cutoffs are $\pm1.658$. A numerical value for $s$ does not supply a known population $\sigma$; the large sample only makes the $t$ cutoffs closer to normal values.
\secInicio[sec:c9-wages][sec:c9-contextual]

\section{Mixed Practice}
\label{sec:c9-mixed}
Choose and justify the distribution before calculating. Use $W$ for the appropriate standardized sample mean. All samples satisfy the introduction's assumptions. Report cutoffs to three decimals, normal areas to four decimals, and probability ranges without interpolation. Solutions are provided after the complete exercise set.
\SectionProblemsTableOfContents{
\SectionProblemLink{sec:c9-mixed-tasks}{Exercise set} \\
\SectionProblemLink{sec:c9-mixed-solutions}{Solutions}
}
\vspace{5mm}

\subsection{Exercise set}
\label{sec:c9-mixed-tasks}
\QuestionsTasks[sec:c9-mixed]{
{9.1}/{sol:c9-m1}/{A retail sales model uses $n=12$ observations and known population $\sigma=15$ monetary units. Find $c$ such that $P(W<c)=0.05$, then calculate $P(-1.00<W<1.50)$.},
{9.2}/{sol:c9-m2}/{A price survey uses $n=12$ shops; unknown population $\sigma$ is estimated by $s=15$ monetary units. Find symmetric cutoffs with total outside area $0.02$ and estimate $P(W<-2.201)$.},
{9.3}/{sol:c9-m3}/{A payroll study selects $5$ employees from each of $5$ departments. Population $\sigma$ is unknown and estimated using the sample. Find $\nu$, the cutoff exceeded in $1\%$ of outcomes, and the central area between $-2.064$ and $2.064$.},
{9.4}/{sol:c9-m4}/{An electricity-expenditure model uses $n=60$ households and a known population standard deviation. Find the symmetric cutoffs enclosing $0.99$ and calculate $P(W<-2.30)$.},
{9.5}/{sol:c9-m5}/{A delivery-cost model uses $n=20$ observations and a sample standard deviation to estimate unknown population variability. Bound $P(W>2.30)$, $P(|W|>2.30)$, and $P(-2.30<W<2.30)$.},
{9.6}/{sol:c9-m6}/{Two studies of mean investment returns each use $n=10$ observations. A knows population $\sigma=2$ percentage points; B estimates unknown $\sigma$ using $s=2$ percentage points. Find the cutoffs enclosing central area $0.95$ for each. At the fixed cutoff $2.00$, calculate or bound each upper-tail area and identify which is larger.},
{9.7}/{sol:c9-m7}/{A business-cost study uses $n=6$ observations with unknown population standard deviation estimated from the sample. Find the negative cutoff leaving $0.005$ below it and estimate $P(-2.015<W<2.571)$.},
{9.8}/{sol:c9-m8}/{A statistical simulation repeatedly samples $n=60$ independent observations from a normal population with unknown standard deviation, estimated separately in each sample. Find cutoffs enclosing the middle $90\%$. Compare their magnitudes with those for a model having known population standard deviation and calculate the difference. Explain whether increasing the sample size to $120$ changes the distribution choice.}
}
\secInicio[sec:c9-mixed-tasks][sec:c9-mixed]

\subsection{Solutions}
\label{sec:c9-mixed-solutions}
\subsubsection*{C9.1 -- Retail sales}
\phantomsection\label{sol:c9-m1}
Known $\sigma$ gives $Z$. The left-tail cutoff is $c\approx-1.645$ (column $q=0.05$). By symmetry, $\Phi(-1.00)=1-0.8413=0.1587$, so the interval area is $0.9332-0.1587=0.7745$.
\secInicio[sec:c9-mixed-tasks][sec:c9-mixed]

\subsubsection*{C9.2 -- Prices}
\phantomsection\label{sol:c9-m2}
Unknown $\sigma$ estimated by $s$ gives $t_{11}$, with $\nu=12-1=11$. Total outside area $0.02$ gives $q=0.01$, so cutoffs are $\pm2.718$. Row $11$, column $0.025$, and symmetry give $P(W<-2.201)\approx0.025$.
\secInicio[sec:c9-mixed-tasks][sec:c9-mixed]

\subsubsection*{C9.3 -- Payroll}
\phantomsection\label{sol:c9-m3}
$n=5\times5=25$ and $\nu=24$. Use $t_{24}$ because population $\sigma$ is estimated. Column $0.01$ gives the positive cutoff $2.492$. The value $2.064$ is in column $0.025$, giving central area approximately $1-2(0.025)=0.95$.
\secInicio[sec:c9-mixed-tasks][sec:c9-mixed]

\subsubsection*{C9.4 -- Electricity expenditure}
\phantomsection\label{sol:c9-m4}
Use $Z$ because $\sigma$ is known. Central area $0.99$ gives $q=0.005$ and cutoffs $\pm2.576$. Symmetry gives $P(W<-2.30)=1-0.9893\approx0.0107$.
\secInicio[sec:c9-mixed-tasks][sec:c9-mixed]

\subsubsection*{C9.5 -- Delivery costs}
\phantomsection\label{sol:c9-m5}
Use $t_{19}$ because $s$ estimates unknown $\sigma$ and $\nu=20-1=19$. The entries $2.093<2.30<2.539$ give
\[
0.01<P(W>2.30)<0.025,\quad
0.02<P(|W|>2.30)<0.05,\quad
0.95<P(-2.30<W<2.30)<0.98.
\]
\secInicio[sec:c9-mixed-tasks][sec:c9-mixed]

\subsubsection*{C9.6 -- Investment comparison}
\phantomsection\label{sol:c9-m6}
A uses $Z$ because $\sigma$ is known; B uses $t_9$ because $s$ estimates $\sigma$ and $\nu=9$. Central area $0.95$ gives $q=0.025$, hence $\pm1.960$ for A and $\pm2.262$ for B. At $2.00$, A's upper-tail area is $1-0.9772=0.0228$. For B, $1.833<2.00<2.262$ gives $0.025<P(W>2.00)<0.05$. B's tail is larger, since even its lower bound exceeds $0.0228$.
\secInicio[sec:c9-mixed-tasks][sec:c9-mixed]

\subsubsection*{C9.7 -- Business costs}
\phantomsection\label{sol:c9-m7}
Use $t_5$ because $\sigma$ is estimated from the sample; $\nu=6-1=5$. The negative cutoff is $-4.032$, using column $0.005$. The asymmetric interval excludes a left tail of $0.05$ and a right tail of $0.025$, giving area approximately $0.925$.
\secInicio[sec:c9-mixed-tasks][sec:c9-mixed]

\subsubsection*{C9.8 -- Statistical simulation}
\phantomsection\label{sol:c9-m8}
Use $t_{59}$ because each sample estimates unknown $\sigma$; $\nu=60-1=59$. Central area $0.90$ gives $q=0.05$, hence cutoffs $\pm1.671$. Known $\sigma$ would give $Z$ cutoffs $\pm1.645$; the magnitude difference is $0.026$. With $n=120$, $\sigma$ remains unknown, so use $t_{119}$ and cutoffs $\pm1.658$. Their magnitude difference from the normal cutoffs is $0.013$: closer, but the choice remains $t$.
\secInicio[sec:c9-mixed-tasks][sec:c9-mixed]

\end{document}
